\sect{माघ-कृष्ण-षट्तिला-एकादशी-माहात्म्यम्}
\label{sec:padma-magha-krishna-shattila}

\ganapatyadiDhyanam
\padmaPuranam

\dnsub{अथ चतुश्चत्वारिंशोऽध्यायः॥४४}

\uvacha{युधिष्ठिर उवाच}

\twolineshloka
{साधु कृष्ण जगन्नाथ आदिदेव जगत्पते}
{कथयस्व प्रसादेन कृपां कुरु ममोपरि}% ॥१॥

\twolineshloka
{माघस्य कृष्णपक्षे तु का वा चैकादशी भवेत्}
{किं नाम को विधिस्तस्या  एतद्विस्तरतो वद}% ॥२॥

\uvacha{श्रीभगवानुवाच}

\twolineshloka
{शृणु त्वं नृपशार्दूल कृष्णमाघस्य या भवेत्}
{षट्तिला नाम विख्याता सर्वपापप्रणाशिनी}% ॥१॥

\twolineshloka
{षट्तिलायाः शृणुष्व त्वं कथां पापहरां शुभाम्}
{यां पुलस्त्यो मुनिश्रेष्ठो दाल्भ्यं प्रति समुक्तवान्}% ॥२॥

\uvacha{दाल्भ्य उवाच}

\twolineshloka
{मर्त्यलोकमनुप्राप्ताः पापं कुर्वन्ति जन्तवः}
{ब्रह्महत्यादिपापैश्च युक्ता ये गोवधादिभिः}% ॥३॥

\twolineshloka
{परद्रव्यापहाराश्च परव्यसनमोहिताः}
{कथं न यान्ति नरकं ब्रह्मंस्तद् ब्रूहि तत्त्वतः}% ॥४॥

\twolineshloka
{अनायासेन भगवन् दानेनाल्पेन केनचित्}
{पापं प्रशमनं याति एतन्मे वक्तुमर्हसि}% ॥५॥

\uvacha{पुलस्त्य उवाच}

\twolineshloka
{साधु साधु महाभाग गुह्यमेतत् सुदुर्लभम्}
{यन्न कस्यचिदाख्यातं विष्णुब्रह्मेन्द्रदैवतैः}% ॥६॥

\twolineshloka
{तदहं कथयिष्यामि त्वया पृष्टो द्विजोत्तम}
{माघमासे तु सम्प्राप्ते शुचिस्नातो जितेन्द्रियः}% ॥७॥

\twolineshloka
{कामक्रोधाभिमानेर्ष्यालोभपैशुन्यवर्जितः}
{देवदेवं च संस्मृत्य पादौ प्रक्षाल्य वारिभिः}% ॥८॥

\twolineshloka
{भूमावपतितं गृह्य गोमयं तत्र मानवः}
{तिलान् प्रक्षिप्य कार्पांसं पिण्डिकाश्चैव कारयेत्}% ॥९॥

\twolineshloka
{अष्टोत्तरशतं चैव नात्र कार्या विचारणा}
{ततो माघे च सम्प्राप्ते ह्याषाढर्क्षं भवेद्यदि}% ॥१०॥

\twolineshloka
{मूलं वा कृष्णपक्षस्यैकादशी नियमांस्ततः}
{गृह्णीयात् पुण्यकाले च विधानं तत्र मे शृणु}% ॥११॥

\twolineshloka
{देवदेवं समभ्यर्च्य सुस्नातः प्रयतः शुचिः}
{कृष्णनामानि सङ्कीर्त्य पुनः प्रस्खलितादिषु}% ॥१२॥

\twolineshloka
{रात्रौ जागरणं कुर्यादादौ होमं च कारयेत्}
{अर्चयेद् देवदेवेशं द्वितीयेऽह्नि पुनर्हरिम्}% ॥१३॥

\twolineshloka
{चन्दनागुरुकर्पूरैर्नैवेद्यं कृसरं तथा}
{संस्मृत्य नाम्ना च ततः कृष्णाख्येन पुनः पुनः}% ॥१४॥

\threelineshloka
{कूष्माण्डैर्नारिकेरैश्च ह्यथवा बीजपूरकैः}
{सर्वाभावेऽपि विप्रेन्द्र शस्तपूगफलैर्वृतम्}
{अर्घं दद्याद्विधानेन पूजयित्वा जनार्दनम्}% ॥१५॥

\twolineshloka
{कृष्ण कृष्ण कृपालुस्त्वमगतीनां गतिर्भव}
{संसारार्णवमग्नानां प्रसीद पुरुषोत्तम}% ॥१६॥

\threelineshloka
{नमस्ते पुण्डरीकाक्ष नमस्ते विश्वभावन}
{सुब्रह्मण्य नमस्तेऽस्तु महापुरुषपूर्वज}
{गृहाणार्घ्यं मया दत्तं लक्ष्म्या सह जगत्पते}% ॥१७॥
[इत्यर्घमन्त्रः]

\twolineshloka
{ततस्तु पूजयेद्विप्रमुदकुम्भं प्रदापयेत्}
{छत्रोपानहवस्त्रैश्च कृष्णो मे प्रीयतामिति}% ॥१८॥

\twolineshloka
{कृष्णा धेनुः प्रदातव्या यथाशक्ति द्विजोत्तमे}
{तिलपात्रं द्विजश्रेष्ठ दद्यात् पात्रविचक्षणः}% ॥१९॥

\twolineshloka
{स्नाने प्राशनके शस्तास्तथा कृष्णतिला मुने}
{तान् प्रदद्यात् प्रयत्नेन यथाशक्ति द्विजोत्तमे}

\twolineshloka
{तिलप्ररोहजाः क्षेत्रे यावत् सङ्ख्यास्तिला द्विज}% ॥२०॥
{तावद् वर्षसहस्राणि स्वर्गलोके महीयते}

\twolineshloka
{तिलस्नायी तिलोद् वर्ती तिलहोमी तिलोदकी}% ॥२१॥
{तिलदाता च भोक्ता च षट्तिला पापनाशिनी}

\uvacha{नारद उवाच}

\onelineshloka
{कृष्ण कृष्ण महाबाहो नमस्ते विश्वभावन}% ॥२२॥

\twolineshloka*
{षट्तिलैकादशीभूतं कीदृशं फलमस्ति वै}
{सोपाख्यानं मम ब्रूहि यदि तुष्टोऽसि यादव}

\uvacha{श्रीकृष्ण उवाच}

\onelineshloka*
{शृणु राजन् यथावृत्तं दृष्टं तत्कथयामि ते}

\twolineshloka
{मर्त्यलोके पुरा ह्यासीद् ब्राह्मण्येका च नारद}
{व्रतचर्यारता नित्यं देवपूजारता सदा}% ॥२३॥

\twolineshloka
{मासोपवासनिरता मम भक्ता च सर्वदा}
{कृष्णोपवाससंयुक्ता मम पूजापरायणा}% ॥२४॥

\twolineshloka
{शरीरं क्लेशितं चैव उपवासैर्द्विजोत्तम}
{देवानां ब्राह्मणानां च कुमारीणां च भक्तितः}% ॥२५॥

\twolineshloka
{गृहादिकं प्रयच्छन्ती सर्वकालं महासती}
{अतिकृच्छ्ररता सा तु सर्वकालं तु वै द्विज}% ॥२६॥

\onelineshloka
{न दत्ता भिक्षुके भिक्षा ब्राह्मणा न च तर्पिताः}% ॥२७॥

\twolineshloka
{ततः कालेन महता मया वै चिन्तितं द्विज}
{शुद्धमस्याः शरीरं हि व्रतैः कृच्छ्रैर्न संशयः}% ॥२८॥

\twolineshloka
{अर्चितो वैष्णवो लोकः कायक्लेशेन  वै तया}
{न दत्तमन्नदानं हि येन तृप्तिः परा भवेत्}% ॥२९॥

\twolineshloka
{एवं ज्ञात्वा अहं ब्रह्मन् मर्त्यलोकमुपागतः}
{कापालं रूपमास्थाय भिक्षापात्रे च याचिता}% ॥३०॥

\twolineshloka
{कस्मात् त्वमागतो ब्रह्मन् क्व यासि वद मेऽग्रतः}
{पुनरेव मया प्रोक्तं देहि भिक्षां च सुन्दरि}% ॥३१॥

\twolineshloka
{तया कोपेन महता मृत्पिण्डस्ताम्रभाजने}
{क्षिप्तो यावदहं ब्रह्मन् पुनः स्वर्गं गतो द्विज}% ॥३२॥

\twolineshloka
{ततः कालेन महता तापसी सुमहाव्रता}
{सदेहा स्वर्गमायाता व्रतचर्याप्रभावतः}% ॥३३॥

\twolineshloka
{मृत्पिण्डिकाप्रदानेन गृहं प्राप्तं मनोरमम्}
{सञ्जातं चैव विप्रर्षे धान्यराशिविवर्जितम्}% ॥३४॥

\twolineshloka
{गृहं यावन्निरीक्षेत न किञ्चित् तत्र पश्यति}
{तावद् गृहाद् विनिष्क्रान्ता ममान्ते चाऽऽगता द्विज}% ॥३५॥

\twolineshloka
{क्रोधेन महताऽऽविष्टमिदं वचनमब्रवीत्}
{मया व्रतैश्च कृच्छ्रैश्च उपवासैरनेकशः}% ॥३६॥

\twolineshloka
{पूजयाऽऽराधितो देवः सर्वलोकस्य पालकः}
{न तत्र  दृश्यते किञ्चिद् गृहे मम जनार्दन}% ॥३७॥

\twolineshloka
{ततश्चोक्तं मया तस्यै गृहं गच्छ महाव्रते}
{आगमिष्यन्ति सुतरां कौतूहलसमन्विताः}% ॥३८॥

\twolineshloka
{देवपत्न्यो हि द्रष्टुं त्वां विस्मयाभिसमन्विताः}
{द्वारं नोद्घाटय विना षट्तिलापुण्यवाचनात्}% ॥३९॥

\twolineshloka
{एवमुक्ता मया सा तु गता वै मानुषी तदा}
{अत्रान्तरे समायाता देवपत्न्यश्च पाण्डव}% ॥४०॥

\twolineshloka
{ताभिश्च कथितं तत्र त्वां द्रष्टुं हि समागताः}
{द्वारमुद्घाटयस्वाद्य त्वां प्रपश्याम शोभने}% ॥४१॥

\uvacha{मानुष्युवाच}

\twolineshloka
{यदि मद्दर्शनं कार्यं सत्यं वाच्यं विशेषतः}
{षट्तिलाया व्रतं पुण्यं द्वारोद्घाटनकारणात्}% ॥४२॥

\uvacha{श्रीकृष्ण उवाच}

\twolineshloka
{एकाऽपि नावदत् तत्र षट्तिलैकादशीव्रतम्}
{अन्यया कथितं तत्र द्रष्टव्या मानुषी मया}% ॥४३॥

\twolineshloka
{ततो द्वारं समुद्घाट्य दृष्टा ताभिश्च मानुषी}
{न देवी न च गन्धर्वी नासुरी न च पन्नगी}% ॥४४॥

\twolineshloka
{दृष्टा पूर्वं तथा नारी यादृशीयं द्विजर्षभ}
{देवीनामुपदेशेन षट्तिलाया व्रतं कृतम्}% ॥४५॥

\twolineshloka
{मानुष्या सत्यव्रतया भुक्तिमुक्तिफलप्रदम्}
{रूपकान्तिसमायुक्ता क्षणेन समवाप सा}% ॥४६॥

\twolineshloka
{धनं धान्यं च वस्त्रादि सुवर्णं रौप्यमेव च}
{भवनं सर्वसम्पन्नं षट्तिलायाः प्रभावतः}% ॥४७॥

\onelineshloka
{रूपकान्तिसमायुक्ता क्षणेन समपद्यत}% ॥४८॥

\twolineshloka
{अतितृष्णा न कर्तव्या वित्तशाठ्यं विवर्जयेत्}
{आत्मवित्तानुसारेण तिलान् वस्त्राणि दापयेत्}% ॥४९॥

\twolineshloka
{लभते चैवमारोग्यं नरो जन्मनि जन्मनि}
{न दारिद्र्यं न कष्टत्वं न च दौर्भाग्यमेव च}% ॥५०॥

\twolineshloka
{सम्भवेद् वै द्विजश्रेष्ठ षट्तिलासमुपोषणात्}
{अनेन विधिना भूप तिलदाता न संशयः}% ॥५१॥

\threelineshloka
{मुच्यते पातकैः सर्वैरनायासेन मानवः}
{दानं च विधिवत् पात्रे सर्वपातकनाशनम्}
{नानर्थः कश्चिन्नाऽऽयासः शरीरे नृपसत्तम}% ॥५२॥


॥इति श्रीपाद्मे महापुराणे पञ्चपञ्चाशत्साहस्र्यां संहितायामुत्तरखण्डे उमापतिनारदसंवादान्तर्गतकृष्णयुधिष्ठिरसंवादे माघ-कृष्ण-षट्तिला-एकादशी-माहात्म्यं नाम चतुश्चत्वारिंशोऽध्यायः॥४४॥

\genMangalaShloka

\hyperref[sec:ekadashi_mahatmyam_padma_puranam]{\closesub}
\clearpage

